\subsection{The abelian approximation in the manuscript}

\begin{statement}{A.17}{The torus quotient and its equivariant splitting}{Proved}
Let $U_n/\Q$ be the plus Fermat curve with its three coordinate divisors
removed, let $F=\pi_1(U_{1,\overline\Q})$, and identify
$N_n=\pi_1(U_{n,\overline\Q})$ with the inverse image of
$n\widehat\Z^2$ under $F\to F^{\mathrm{ab}}$.
The quotient of the arithmetic fundamental extension by the closed
subgroup $F'\subset N_n$ is canonically the arithmetic fundamental
extension of
\[
 T=\mathbb G_m^{\{0,1,\infty\}}/\mathbb G_m.
\]
Consequently its geometric-conjugacy classes of sections over $G_\Q$
are a torsor, canonically based by $[1:1:1]$, under
$H^1(\Q,\widehat\Z(1))^2=(\widehat{\Q^\times})^2$.
After retaining the permutation action of $\mathfrak S_3$, the extension
over $G_\Q\times\mathfrak S_3$ splits, and any two continuous splittings
are conjugate by a unique geometric element. These assertions verify R8
for the quotient in that statement; they do not assert a splitting of
the full Fermat extension.
\end{statement}
\begin{proof}
Embed $U_n$ into $T$ by its three nonzero homogeneous coordinates.
The square formed with $U_n\to U_1$ and the power map
$[n]:T\to T$ is cartesian, in the category of schemes over $\Q$.
Over $\overline\Q$, the two coordinate ratios identify the abelianization
of the tripod group with the geometric fundamental group of $T$.
Indeed their winding numbers at the coordinate divisors generate the
rank-two lattice with its sole relation that the three peripheral
vectors sum to zero. It follows that the kernel of
$N_n\to\pi_1(T_{\overline\Q})$ is exactly $F'$, and the map is
surjective: multiplication by $n$ identifies the target with
$nF^{\mathrm{ab}}$ inside $F^{\mathrm{ab}}$.
The induced arithmetic homomorphism is also surjective, since it is
surjective on both the geometric kernel and the quotient $G_\Q$.
The same comparison is equivariant for permutations of the coordinates.

The point $[1:1:1]\in T(\Q)$ is fixed by every permutation, so its
decomposition group supplies a splitting over
$G_\Q\times\mathfrak S_3$. Write
\[
 L=\Z^3/\Z(1,1,1),\qquad V=\widehat\Z(1)\otimes_\Z L.
\]
The arithmetic action on $V$ is cyclotomic and the permutation action
is the indicated action on $L$; the two actions commute. Kummer theory,
including passage to the inverse limit of the finite Kummer sequences,
gives $H^1(G_\Q,V)=M\otimes L$, where
$M=\varprojlim_m\Q^\times/(\Q^\times)^m$.
Here no extra $\varprojlim^1$ term occurs in degree one: the finite
groups $H^0(G_\Q,\mu_m)$ form a Mittag--Leffler system along a divisible
cofinal sequence of $m$.

We have $V^{G_\Q}=0$. Inflation--restriction therefore gives
\[
 H^1(G_\Q\times\mathfrak S_3,V)
 \simeq H^1(G_\Q,V)^{\mathfrak S_3}
 \simeq (M^3/\operatorname{diag}M)^{\mathfrak S_3}=0.
\]
For the last equality, an invariant class represented by $(a,b,c)$
has $(b-a,a-b,0)$ diagonal. Its third coordinate forces that diagonal
element to be zero, hence $a=b$; a second transposition gives $a=c$.
This argument works even when $M$ has torsion. Thus every splitting
is conjugate to the displayed one. Its geometric centralizer is
$V^{G_\Q\times\mathfrak S_3}=0$, which gives uniqueness of the
conjugating element. Restricting to $G_\Q$ alone leaves the stated
two-coordinate Kummer torsor. This proves all comparisons, including
their direction: the torus quotient admits many sections although
the full plus curve for even $n$ admits none.
\end{proof}
